% Chapter 19: PDE Solvers
% ch19_pdes.tex

\chapter{Numerical PDE Solvers: Heat, Wave, and Laplace Equations}

\begin{historical}
Euler's \emph{Institutiones Calculi Integralis} (1768--1770) laid the groundwork for solving
differential equations numerically. His finite difference approach to the calculus of variations
anticipated modern PDE solvers by over a century. Today, the same finite difference stencils
Euler pioneered are used to simulate heat transfer, acoustic waves, and electrostatic fields.

Euler's key insight was to replace derivatives with difference quotients on a discrete grid.
For example, the second derivative can be approximated as:
\[
\frac{d^2u}{dx^2} \approx \frac{u(x+h) - 2u(x) + u(x-h)}{h^2}
\]
This simple idea, combined with Taylor series analysis, leads to the modern theory of
numerical PDEs.
\end{historical}

\section{Prerequisites}

To work through this chapter, you need:
\begin{itemize}
    \item \textbf{Mathematics:} Partial derivatives, the heat equation, wave equation, and Laplace equation.
    Understanding of Taylor series for deriving finite difference approximations.
    \item \textbf{Programming:} Python lists for 2D grids, nested loops for stencil operations,
    and basic numerical stability analysis.
\end{itemize}

\section{Derivation of Finite Difference Stencils}

The finite difference method replaces derivatives with algebraic differences on a grid.
Consider a uniform grid with spacing $\Delta x$ and $\Delta y$. For a function $u(x,y)$,
we define grid values $u_{i,j} \approx u(x_i, y_j)$ where $x_i = i\Delta x$ and $y_j = j\Delta y$.

\subsection{Second Derivative Approximation}

Using Taylor series expansions:
\[
u_{i+1,j} = u_{i,j} + \Delta x \frac{\partial u}{\partial x}\bigg|_{i,j} + \frac{\Delta x^2}{2} \frac{\partial^2 u}{\partial x^2}\bigg|_{i,j} + \frac{\Delta x^3}{6} \frac{\partial^3 u}{\partial x^3}\bigg|_{i,j} + \cdots
\]
\[
u_{i-1,j} = u_{i,j} - \Delta x \frac{\partial u}{\partial x}\bigg|_{i,j} + \frac{\Delta x^2}{2} \frac{\partial^2 u}{\partial x^2}\bigg|_{i,j} - \frac{\Delta x^3}{6} \frac{\partial^3 u}{\partial x^3}\bigg|_{i,j} + \cdots
\]

Adding these equations eliminates the odd derivatives:
\[
u_{i+1,j} + u_{i-1,j} = 2u_{i,j} + \Delta x^2 \frac{\partial^2 u}{\partial x^2}\bigg|_{i,j} + O(\Delta x^4)
\]

Therefore:
\[
\frac{\partial^2 u}{\partial x^2}\bigg|_{i,j} \approx \frac{u_{i+1,j} - 2u_{i,j} + u_{i-1,j}}{\Delta x^2}
\]

This is a \textbf{second-order accurate} approximation (error $O(\Delta x^2)$).

\section{The Heat Equation}

The 2D heat equation describes temperature diffusion:
\[
\frac{\partial u}{\partial t} = \alpha \left( \frac{\partial^2 u}{\partial x^2} + \frac{\partial^2 u}{\partial y^2} \right)
\]
where $\alpha$ is thermal diffusivity. Euler's finite difference approach replaces derivatives
with discrete approximations:
\[
\frac{u_i^{n+1} - u_i^n}{\Delta t} = \alpha \frac{u_{i+1}^n - 2u_i^n + u_{i-1}^n}{\Delta x^2} + \alpha \frac{u_{i,j+1}^n - 2u_{i,j}^n + u_{i,j-1}^n}{\Delta y^2}
\]

This yields the explicit FTCS (Forward Time, Central Space) scheme:
\[
u_i^{n+1} = u_i^n + r_x(u_{i+1}^n - 2u_i^n + u_{i-1}^n) + r_y(u_{i,j+1}^n - 2u_{i,j}^n + u_{i,j-1}^n)
\]
where $r_x = \alpha \Delta t / \Delta x^2$ and $r_y = \alpha \Delta t / \Delta y^2$.

\begin{lstlisting}[language=Python, caption=Explicit heat equation solver]
from euler import heat_equation_explicit
import math

# Initial condition: hot spot in center
nx, ny = 40, 40
u0 = [[0.0]*nx for _ in range(ny)]
u0[20][20] = 100.0  # hot spot

# Solve with alpha=0.1, dx=0.1, dt=0.005
solutions = heat_equation_explicit(
    u0, alpha=0.1, dx=0.1, dt=0.005,
    nx=40, ny=40, n_steps=200,
    bc_left=0.0, bc_right=0.0,
    bc_top=0.0, bc_bottom=0.0
)

print(f"Generated {len(solutions)} time steps")
print(f"Peak temp at t=0: {max(max(r) for r in solutions[0]):.1f}")
print(f"Peak temp at t=final: {max(max(r) for r in solutions[-1]):.4f}")
\end{lstlisting}

\textbf{Stability:} The explicit scheme requires $\alpha \Delta t / \Delta x^2 \leq 0.25$ in 2D.
Violating this produces oscillations that grow exponentially. This condition arises from
requiring the coefficients in the update formula to be non-negative (maximum principle).

\subsection{Stability Analysis}

For the 1D heat equation, von Neumann analysis gives the amplification factor:
\[
G = 1 - 4r \sin^2\left(\frac{k \Delta x}{2}\right)
\]
where $r = \alpha \Delta t / \Delta x^2$. Stability requires $|G| \leq 1$, which gives:
\[
-1 \leq 1 - 4r \sin^2\left(\frac{k \Delta x}{2}\right) \leq 1
\]
The right inequality is always satisfied. The left inequality requires:
\[
4r \sin^2\left(\frac{k \Delta x}{2}\right) \leq 2
\]
The worst case is $\sin^2(k\Delta x/2) = 1$, giving $r \leq 0.5$ for 1D and $r \leq 0.25$ for 2D.

\section{Crank--Nicolson Method}

The Crank--Nicolson method averages the explicit and implicit schemes, yielding second-order
accuracy in time and unconditional stability:
\[
\frac{u^{n+1} - u^n}{\Delta t} = \frac{\alpha}{2} \left( \nabla^2 u^{n+1} + \nabla^2 u^n \right)
\]

This requires solving a linear system at each step. The ADI (Alternating Direction Implicit)
splitting makes this efficient by solving one-dimensional systems alternately.

\begin{lstlisting}[language=Python, caption=Crank-Nicolson solver]
from euler import heat_equation_crank_nicolson

# Same initial condition as before
solutions = heat_equation_crank_nicolson(
    u0, alpha=0.1, dx=0.1, dt=0.01,
    nx=40, ny=40, n_steps=500
)
# No stability restriction! Can use larger dt.
\end{lstlisting}

\textbf{Complexity:} Each Crank--Nicolson step requires solving a tridiagonal system in each
direction. Using the Thomas algorithm, this takes $O(N)$ operations per direction, where $N$
is the number of grid points. For a $200 \times 200$ grid, each time step requires approximately
$80{,}000$ operations.

\section{The Wave Equation}

The 2D wave equation describes vibrating membranes and sound propagation:
\[
\frac{\partial^2 u}{\partial t^2} = c^2 \left( \frac{\partial^2 u}{\partial x^2} + \frac{\partial^2 u}{\partial y^2} \right)
\]

Using central differences in both time and space:
\[
\frac{u_i^{n+1} - 2u_i^n + u_i^{n-1}}{\Delta t^2} = c^2 \frac{u_{i+1}^n - 2u_i^n + u_{i-1}^n}{\Delta x^2} + c^2 \frac{u_{i,j+1}^n - 2u_{i,j}^n + u_{i,j-1}^n}{\Delta y^2}
\]

\begin{lstlisting}[language=Python, caption=Wave equation solver]
from euler import wave_equation_explicit
import math

# Initial displacement: Gaussian bump
nx, ny = 60, 60
u0 = [[math.exp(-((i-30)**2 + (j-30)**2)/200.0)
       for j in range(nx)] for i in range(ny)]
v0 = [[0.0]*nx for _ in range(ny)]  # zero initial velocity

solutions = wave_equation_explicit(
    u0, v0, c=1.0, dx=0.1, dt=0.05,
    nx=60, ny=60, n_steps=200,
    bc_left=0.0, bc_right=0.0,
    bc_top=0.0, bc_bottom=0.0
)
\end{lstlisting}

\textbf{CFL condition:} For 2D stability, $c \Delta t / \Delta x \leq 1/\sqrt{2}$.
This is the Courant--Friedrichs--Lewy condition, which requires that information
cannot travel more than one grid cell per time step.

\section{Laplace and Poisson Equations}

The Laplace equation $\nabla^2 u = 0$ describes steady-state temperature, electrostatic
potential, and incompressible flow. The Jacobi iteration updates each interior point to
the average of its neighbors:
\[
u_i^{(k+1)} = \frac{1}{4} \left( u_{i+1}^{(k)} + u_{i-1}^{(k)} + u_{i,j+1}^{(k)} + u_{i,j-1}^{(k)} \right)
\]

\begin{lstlisting}[language=Python, caption=Laplace equation solver]
from euler import laplace_equation_jacobi

# 20x20 grid, top edge at temperature 1, rest at 0
u0 = [[0.0]*20 for _ in range(20)]
u0[0] = [1.0]*20  # top boundary

sol, iterations = laplace_equation_jacobi(
    u0, nx=20, ny=20,
    bc_top=1.0, bc_bottom=0.0,
    bc_left=0.0, bc_right=0.0,
    tol=1e-8
)
print(f"Converged in {iterations} iterations")
print(f"Center temperature: {sol[10][10]:.4f}")
\end{lstlisting}

\subsection{Convergence of Jacobi Iteration}

Jacobi iteration converges if the spectral radius of the iteration matrix is less than 1.
For the Laplace equation on an $N \times N$ grid, the convergence rate is:
\[
\rho \approx 1 - \frac{\pi^2}{2N^2}
\]
This means the error reduces by a factor of approximately $1 - O(1/N^2)$ per iteration.
For $N = 100$, we need roughly $5000$ iterations to reduce the error by a factor of $e$.

\subsection{Gauss--Seidel and SOR}

The \textbf{Gauss--Seidel} method uses newly computed values immediately:
\[
u_i^{(k+1)} = \frac{1}{4} \left( u_{i+1}^{(k)} + u_{i-1}^{(k+1)} + u_{i,j+1}^{(k)} + u_{i,j-1}^{(k+1)} \right)
\]

This typically converges about twice as fast as Jacobi. \textbf{Successive Over-Relaxation (SOR)}
accelerates convergence further:
\[
u_i^{(k+1)} = (1-\omega)u_i^{(k)} + \omega \cdot u_i^{\text{GS}}
\]
With optimal $\omega \approx 2/(1 + \pi/N)$, SOR converges orders of magnitude faster than Jacobi.

\section{Von Neumann Stability Analysis}

Von Neumann analysis examines the growth of Fourier modes to determine stability.
For the explicit heat equation, the amplification factor is:
\[
G = 1 - 4r \sin^2\left(\frac{k_x \Delta x}{2}\right) - 4r \sin^2\left(\frac{k_y \Delta y}{2}\right)
\]
where $r = \alpha \Delta t / \Delta x^2$. Stability requires $|G| \leq 1$ for all modes.

\begin{lstlisting}[language=Python, caption=Stability analysis]
from euler import von_neumann_stability

# Check heat equation stability
info = von_neumann_stability('heat_explicit',
    alpha=0.1, dx=0.1, dt=0.01)
print(f"Stable: {info['stable']}")
print(f"Constraint: {info['constraints'][0]}")

# Check wave equation stability
info2 = von_neumann_stability('wave_explicit',
    c=1.0, dx=0.1, dt=0.05)
print(f"Courant number: {info2.get('courant', 'N/A')}")
\end{lstlisting}

\section{1D Solvers}

For simpler problems, 1D solvers are faster and easier to analyze:

\begin{lstlisting}[language=Python, caption=1D heat and wave solvers]
from euler import solve_heat_1d, solve_wave_1d
import math

# 1D heat: initial parabolic profile, zero BCs
def u0(x): return x * (1.0 - x)
x, t, u = solve_heat_1d(u0, alpha=0.1, L=1.0,
                         T=0.1, nx=100, nt=500)
print(f"Final peak temperature: {max(u):.4f}")

# 1D wave: plucked string
def u0w(x): return math.sin(math.pi * x)
def v0w(x): return 0.0
x, t, u = solve_wave_1d(u0w, v0w, c=1.0, L=1.0,
                          T=2.0, nx=100, nt=500)
print(f"String displacement at t=2: {max(abs(v) for v in u):.4f}")
\end{lstlisting}

\section{Boundary Conditions}

\textbf{Dirichlet conditions} specify the value of $u$ on the boundary (e.g., fixed temperature).
\textbf{Neumann conditions} specify the derivative $\partial u/\partial n$ (e.g., insulated boundary
with zero heat flux).

For Neumann conditions, we introduce \textbf{ghost points} outside the domain and use the
boundary condition to relate them to interior points. For example, a zero-flux boundary at $x=0$:
\[
\frac{u_{1,j} - u_{-1,j}}{2\Delta x} = 0 \implies u_{-1,j} = u_{1,j}
\]

\section{Modern Relevance}

PDE solvers are the backbone of modern engineering simulation. The heat equation models thermal
management in electronics, the wave equation governs seismic analysis and acoustic design, and
Laplace/Poisson solvers are essential for electromagnetic field simulation, structural mechanics,
and computational fluid dynamics. The finite difference methods Euler pioneered remain the simplest
and most intuitive approach, taught in every introductory PDE course worldwide.

\section{Frequently Asked Questions}

\begin{itemize}
\item \textbf{Q: Why does the explicit heat equation have a stability limit?} \\
The explicit scheme updates each point using only the current values of its neighbors.
If $\Delta t$ is too large, errors amplify faster than diffusion can smooth them. The condition
$\alpha \Delta t / \Delta x^2 \leq 0.25$ ensures the update is a weighted average (non-negative
coefficients), preventing oscillations.

\item \textbf{Q: What is the CFL condition?} \\
The Courant--Friedrichs--Lewy condition states that the numerical domain of dependence must
contain the physical domain of dependence. For the wave equation, information travels at speed
$c$, so $\Delta t$ must be small enough that a wave doesn't cross more than one grid cell per
time step: $c \Delta t / \Delta x \leq 1/\sqrt{2}$ in 2D.

\item \textbf{Q: Why is Crank--Nicolson unconditionally stable?} \\
Crank--Nicolson averages the forward and backward Euler schemes. The backward Euler part is
implicit and damps all high-frequency modes. The averaging preserves second-order accuracy while
inheriting the damping property. Mathematically, the amplification factor is always $\leq 1$
regardless of $\Delta t$.

\item \textbf{Q: How does Jacobi iteration converge?} \\
Jacobi iteration is a stationary iterative method that converges if the spectral radius of the
iteration matrix is less than 1. For the Laplace equation on a rectangular grid, convergence is
guaranteed but slow: the error reduces by roughly a factor of $1 - O(1/N^2)$ per iteration, where
$N$ is the number of grid points.

\item \textbf{Q: What is the difference between Dirichlet and Neumann boundary conditions?} \\
\textbf{Dirichlet} specifies the value of $u$ on the boundary (e.g., fixed temperature).
\textbf{Neumann} specifies the derivative $\partial u/\partial n$ (e.g., insulated boundary with
zero heat flux). Our examples use Dirichlet conditions; Neumann conditions require modifying the
stencil near the boundary.

\item \textbf{Q: Why use Gauss--Seidel or SOR instead of Jacobi?} \\
Gauss--Seidel uses the most recently computed values immediately, typically converging about
twice as fast as Jacobi. Successive Over-Relaxation (SOR) accelerates this further by extrapolating:
$u_{new} = (1-\omega)u_{old} + \omega u_{GS}$. With optimal $\omega \approx 2/(1+\pi/N)$, SOR
converges orders of magnitude faster.

\item \textbf{Q: What is the method of lines?} \\
The method of lines discretizes only the spatial derivatives, leaving the time derivative continuous.
This converts the PDE into a system of ODEs, which can then be solved using standard ODE methods
(RK4, implicit methods, etc.). It combines the flexibility of ODE solvers with the accuracy of
spatial discretization.
\end{itemize}

\begin{exercise}
Implement the explicit heat equation solver from scratch (without importing) using a 2D list for
the grid. Verify that $r > 0.25$ produces instability.
\end{exercise}

\begin{exercise}
Solve the 1D wave equation with initial condition $u(x,0) = \sin(\pi x)$ and zero boundary
conditions. Show that after one period $T = 2$, the solution returns to its initial state.
\end{exercise}

\begin{exercise}
Use von Neumann analysis to derive the stability condition for the 1D heat equation. Show that
$r = \alpha \Delta t / \Delta x^2 \leq 0.5$ is required.
\end{exercise}

\begin{exercise}
Solve the 2D Laplace equation on a square domain with boundary conditions $u=1$ on the top edge
and $u=0$ on the other three edges. Compare the convergence rates of Jacobi and Gauss--Seidel.
\end{exercise}

\begin{exercise}
Implement the Crank--Nicolson method for the 1D heat equation using the Thomas algorithm for
tridiagonal systems. Verify second-order accuracy in time by comparing solutions with different
time step sizes.
\end{exercise}

\begin{exercise}
For the wave equation, derive the discrete energy:
\[
E^n = \frac{1}{2}\sum_{i,j} \left[\left(\frac{u_{i,j}^{n+1} - u_{i,j}^n}{\Delta t}\right)^2
+ c^2\left(\frac{(u_{i+1,j}^n - u_{i,j}^n)^2}{\Delta x^2} + \frac{(u_{i,j+1}^n - u_{i,j}^n)^2}{\Delta y^2}\right)\right]
\]
Show that this energy is conserved (up to roundoff) for the explicit scheme.
\end{exercise}

\begin{exercise}
Implement a Neumann boundary condition for the heat equation (insulated boundary).
Show that the total heat energy is conserved when all boundaries are insulated.
\end{exercise}

\begin{exercise}
Compare the convergence rate of Jacobi iteration versus SOR for the Laplace equation
on a $100 \times 100$ grid. What is the optimal relaxation parameter $\omega$?
\end{exercise}